\SourceLeaf{321}
\section*{BEMERKUNGEN.}
\pdfbookmark[1]{Bemerkungen zur geometrischen Seite der ternären Formen}{ternary-bemerkungen}
\BookRule{.15\textwidth}
Neben den vorstehenden Notizen, welche die in der Anzeige von \textsc{Seeber}'s Untersuchungen der ternären Formen gegebenen Gesichtspunkte theilweise weiter entwickeln, sind in der Handschrift mehre eigne mit einem achtzölligen \textsc{Reichenbach}'schen Theodolithen ausgeführte Crystallmessungen aufgezeichnet. Die einzelnen Protokolle enthalten das jedesmalige Datum der Beobachtung, woraus zu ersehen ist, dass diese Untersuchung dem Monat Juli $1831$ angehört.

Aus der Theorie der indifferenten ternären quadratischen Formen findet sich im handschriftlichen Nachlass nur der folgende, wahrscheinlich in der Zeit der Ausarbeitung der Disqu. Arr. aufgezeichnete Lehrsatz
\begin{quote}
‘Omnes transformationes formae ternariae
\[
\left(\begin{smallmatrix}1,&1,&-1\\0,&0,&0\end{smallmatrix}\right)
\]
in se ipsam exhibentur per formulam
\[
\begin{array}{ccc}
\alpha\delta+\beta\gamma&\alpha\beta-\gamma\delta&\alpha\beta+\gamma\delta\\[.5ex]
\alpha\gamma-\beta\delta&\frac12(\alpha\alpha+\delta\delta-\beta\beta-\gamma\gamma)&\frac12(\alpha\alpha+\gamma\gamma-\beta\beta-\delta\delta)\\[.5ex]
\alpha\gamma+\beta\delta&\frac12(\alpha\alpha+\beta\beta-\gamma\gamma-\delta\delta)&\frac12(\alpha\alpha+\beta\beta+\gamma\gamma+\delta\delta)
\end{array}
\]
acceptis $\alpha,\beta,\gamma,\delta$ ita ut fiat $\alpha\delta-\beta\gamma=1$.’
\end{quote}

Es entstehen nemlich alle Transformationen, in denen die neun Coëfficienten ganze Zahlen sind, wenn für $\alpha,\beta,\gamma,\delta$ sowohl alle die der Bedingungsgleichung genügenden ganzen Zahlen und zwar zwei gerade und zwei ungerade gesetzt werden, als auch alle die ungeraden Vielfache von $\surd\frac12$, welche dieselbe Bedingungsgleichung $\alpha\delta-\beta\gamma=1$ erfüllen.

\SourceLeaf{322}
Zu Seite $309$. Chaux carbonatée équiaxe, inverse, contrastante und mixte sind die von \textsc{Hauy} (Traité de Minéralogie $1801$ Tome II pag. $132$, $137$) gebrauchten Benennungen.

Die Tafel der Transformationen der Form $\left(\begin{smallmatrix}5,&5,&5\\1,&1,&1\end{smallmatrix}\right)$ enthält in der ersten Verticalreihe die Coëfficienten der Substitution, in der zweiten die dadurch entstandene neue Form, in der dritten die der letztern Form entsprechende primitive, wenn diese nicht selbst schon eine solche ist, und in der vierten deren Adjuncta.
\begin{flushright}\textsc{Schering.}\end{flushright}
\BookRule{.30\textwidth}

